\chapter*{Metode Pengukuran}
\addcontentsline{toc}{chapter}{Metode Pengukuran}

Seluruh angka di buku ini berasal dari program yang disertakan, dijalankan
di satu komputer, dan dicatat apa adanya. Bagian ini menjelaskan cara angka
itu diambil, agar pembaca dapat menilai batas keberlakuannya dan
mengulanginya sendiri.

\section*{Lingkungan pengukuran}

Pengukuran dijalankan di satu komputer meja dengan sistem operasi Linux,
memakai perangkat berikut. Versi yang dipakai pada tiap bab dicatat di bab
itu bila angkanya bergantung pada versinya.

\begin{table}[!htb]
  \centering
  \caption{Perangkat yang dipakai menjalankan pengukuran di buku ini}
  \label{tab:lingkungan}
  \begin{tabularx}{\textwidth}{Xl}
    \hline
    \textbf{Perangkat} & \textbf{Versi} \\
    \hline
    Python, untuk aplikasi contoh dan alat ukur & 3.14 \\
    \hline
    PostgreSQL, untuk bab basis data & 18.1 \\
    \hline
    Docker, untuk bab \textit{deployment} & 29.9 \\
    \hline
    Google Chrome, untuk pengukuran sisi browser & 157 \\
    \hline
    Node.js beserta TypeScript, untuk pemeriksaan tipe & 20.20 dan 7.0 \\
    \hline
  \end{tabularx}
\end{table}

Aplikasi contoh tiap bab berjalan di \texttt{127.0.0.1} pada portnya
sendiri, dan alat ukurnya diarahkan ke alamat itu. Tidak ada pengukuran yang
diarahkan ke sistem milik orang lain.

\section*{Aturan yang dipegang}

Empat aturan dipegang sepanjang buku, dan aturan itulah yang membedakan
angka di sini dari angka yang dikira-kira.

\begin{enumerate}
  \item Diukur lebih dahulu, baru disimpulkan. Tanpa pengukuran, bagian yang
        diperbaiki biasanya bagian yang paling mudah ditebak, bukan bagian
        yang paling mahal.
  \item Keluaran program dikutip sebagaimana tercetak. Bila nama di dalam
        kode berubah dan nama itu ikut tercetak, programnya dijalankan ulang
        dan angka di naskah diperbarui dari hasil yang baru.
  \item Hasil yang janggal tidak dibuang. Angka yang melawan harapan tetap
        dicetak, lalu sebabnya dijelaskan di paragraf sesudahnya.
  \item Kode di naskah adalah salinan dari file yang dijalankan. Potongan
        yang ditampilkan boleh lebih pendek daripada filenya, tetapi tidak
        boleh berbeda isinya.
\end{enumerate}

\section*{Batas keberlakuan angkanya}

Angka pada buku ini terikat pada satu komputer dan satu beban. Tiga hal
berikut perlu diingat ketika angkanya dibandingkan dengan hasil sendiri.

Pertama, nilai milidetik akan berbeda di komputer lain, dan beberapa kali
lebih besar pada telepon kelas bawah. Yang dapat dipakai ulang adalah pola
perbandingannya, misalnya bahwa satu varian tumbuh mengikuti besar halaman
sedangkan varian lain tidak.

Kedua, sebagian pengukuran memuat beban yang sengaja diperlambat atau
kegagalan yang sengaja dibuat, agar akibatnya terlihat dalam satu sesi
latihan. Tiap kali hal itu dilakukan, naskahnya menyebutkannya di dekat
angkanya.

Ketiga, angka jaringan diambil di \texttt{localhost}, sehingga tidak memuat
waktu tempuh antarkota. Bab yang membahas jarak menyebutkan asumsinya
sendiri, dan menghitung akibatnya dengan angka yang dinyatakan terbuka.

\section*{Cara mengulangi pengukurannya}

Seluruh program tersimpan per bab di dalam \berkas{sources/}, dan tiap
direktori punya \berkas{README.md} berisi daftar filenya beserta perintah
untuk menjalankannya. Skrip Python dijalankan memakai \textit{virtual
environment} di \berkas{sources/.venv}. Prosedur lengkap tiap bab, beserta
lembar isian untuk mencatat hasil sendiri, ada pada lampiran di belakang
buku ini.
